\section{Background}
\label{sec:background}

Before diving into the technical details of \mango, we provide the necessary background information on vulnerabilities in binary firmware services and discuss the associated challenges in discovering them with state-of-the-art static analysis techniques.

\subsection{Firmware Types}

Qasem et al.~\cite{qasem2021automatic} categorized firmware into three types:
Type~I, where there is a clear separation between the Operating System (OS) and the contained firmware services (in terms of files and processes), Type~II, where the OS and services are combined into a monolithic binary blob, and Type~III, where there is no OS but only logic embedded in a blob.

Type~II and Type~III firmware are typically used to implement software for programmable logic controllers (PLCs) and real-time, low-power embedded devices (generally with less complicated logic).
Type~I firmware is more commonly seen in feature-rich and complex IoT devices such as routers, webcams, or networked CCTVs, where the OS is usually embedded Linux and services are ELF binaries (of varying architectures).
A 2014 study by Costin et al.~\cite{costin2014large} found that over 86\% of firmware samples in a large dataset of 32,356 images were Type~I firmware using embedded Linux as the OS.
Similar to prior work~\cite{celik2018sensitive,mandal2020cross} in this area, our research focuses on finding taint-style vulnerabilities in the user-space services of Type~I firmware.


\subsection{Motivating Example}
\label{sec:motivating_example}

\begin{code}
    \caption{Decompiled code of the \texttt{httpd} and \texttt{dlnad} binary from the Netgear R6400. \mango can identify this vulnerability while \karonte and \satc will not analyze \texttt{dlnad}.}
    \label{dlnad_vuln}
    \inputminted[
        frame=lines,
        fontsize=\footnotesize,
        linenos,
        breaklines
    ]{c}{assets/code/dlnad_vuln.c}
\end{code}


Listing~\ref{dlnad_vuln} shows a vulnerability in the Netgear R6400 router, which is in the dataset for \karonte and the ablation dataset \satc used.
The vulnerability here is straightforward: a command injection based on the \texttt{iserver\_passcode} front-end keyword.
\satc missed this vulnerability, and \karonte would have as well if it supported command injections.
Both of these tools missed this vulnerability because the border binary analysis for both tools does not include \texttt{dlnad} as a potential target.
For example, \satc prioritizes binaries with more references to frontend keywords and limits border binaries to three per firmware image.
Due to the simple functionality and small size of \texttt{dlnad}, it is excluded by the border binary selection algorithm in \satc.

\subsection{Firmware Vulnerability Analysis}
\label{sec:observation}

For finding taint-style vulnerabilities in firmware, prior work proposes two genres of techniques:
(1) firmware rehosting and dynamic analysis and (2) static analysis.

\smallskip
\noindent
\textbf{Firmware rehosting and dynamic analysis.}
Given the limited computing capabilities of IoT devices, researchers attempt to run firmware services on commodity systems, usually using an emulator such as Qemu~\cite{qemu} to leverage dynamic analysis techniques such as fuzzing.
% Current dynamic vulnerability discovery techniques cannot detect these embedded vulnerabilities effectively.
While researchers have focused on making taint-style analysis more approachable~\cite{clause2007dytan,ming2016straighttaint,ming2015taintpipe}, the performance of dynamic techniques is limited by the effectiveness of firmware re-hosting, the quality of input seeds, and efficient coverage exploration, making it difficult for dynamic techniques to find taint-style vulnerabilities with reasonable scalability.

Consider the motivating example in Listing~\ref{dlnad_vuln}: to use a dynamic analysis technique, an analyst must first identify the connection that \texttt{dlnad} has to \texttt{httpd} through NVRAM variables.
It would be an unlikely target because it does not actively process user input like a web server would.
Further, even if an analyst wanted to analyze \texttt{dlnad}, it requires a specific combination of NVRAM values and properly formatted user input for the variables to trigger.

\smallskip
\noindent
\textbf{Static analysis.}
Static analyses sacrifice precision for higher scalability.
Analyzing binaries statically requires a deep understanding of the binary format, the underlying architecture, and the operating system for which the binary is built.
A typical static analysis workflow includes disassembling, lifting to an intermediate language (IL), CFG recovery, and various types of advanced analysis atop.

Symbolic execution on binaries~\cite{king1976symbolic,wang2017angr} allows tracking data flows in a binary:
It simulates the execution of the target binary on a value domain of both concrete values and symbolic values (e.g., variables with unknown values), collects path constraints, and solves path constraints to derive possible values for symbolic values.
The caveat is that symbolic execution is computationally expensive and suffers from path explosion~\cite{krishnamoorthy2010tackling}.
Another technique for tracking data flows in a binary is static taint analysis~\cite{saxena2008efficient}.
It marks the source of the data and tracks it through the program, approximating the effects of function calls until reaching sinks.
Static taint analysis may suffer from over-tainting, marking values as tainted when they are not, and under-tainting, missing tainted values.

Depending on the analysis goals, these static analyses may adopt value domains (e.g., the concrete domain where only concrete values are considered, or the value-set domain~\cite{balakrishnan2005codesurfer} where every value is modeled using its bits, lower- and upper-bounds, and its stride) with varied precision.

\smallskip
\noindent
\textbf{Firmware taint analysis.}
State-of-the-art static analysis techniques for finding firmware vulnerabilities, such as \karonte~\cite{karonte} and \satc~\cite{SaTC}, start their analysis from \emph{border binaries}, binaries that they determine to be reachable from the local network.
\satc takes it one step further:
It identifies common keywords in web API interfaces and ignores all other interfaces.

While sometimes users knowingly communicate with their IoT devices via front-end web services (e.g., logging into a router's management portal), this is not the only available input vector for firmware services.
Many services, such as \texttt{miniupnpd}, communicate without frontend interfaces but are open to receiving data from the local network, if not the Internet.
By only focusing on border binaries and common web-API-style endpoints, \karonte and \satc are likely to miss a large portion of vulnerability sources, which leads to false negatives (i.e., missed vulnerabilities).

%However, the keyword approach they take is done in an effort to reduce false positives.
%To cover this command-injection case would require static analysis to be fast enough to be run on many binaries in the firmware without losing precision.


\smallskip

While static analysis is the more viable approach for finding taint-style vulnerabilities, both \karonte and \satc severely limit the potential vulnerability sources (by only analyzing border binaries and input locations with specific keywords), which impacts their ability to find vulnerabilities.
Because \rda is more scalable, \mango takes a different approach:
It analyzes \emph{every valid binary} in each firmware sample, flagging and filtering potential vulnerabilities.

\subsection{Threat Model and Scope}
\label{sec:threat_model}

Unlike \karonte or \satc, we do not limit which binaries handle user input.
Both use the concept of Border Binaries to limit the scope of their analysis to both reduce false positives and analysis times.
However, \mango does not need this limit as it is able to analyze entire firmware samples in a reasonable amount of time.
We do this exclusively on Linux-based user-space firmware binaries.
% We focus on vulnerabilities that an attacker can trigger by supplying input through the following manner: sockets, environment variables, or NVRAM variables with no specific link to front-end files.
% Therefore, while we consider attackers who can only reach the target IoT device over the network, unlike \karonte or \satc, we do not limit \emph{how} attackers reach the target device.
%We focus on any binary that could potentially process external data in a harmful manner, be it through sockets, arguments passed through command lines, or environment and NVRAM variables.


%% To serve users' requests, IoT services must exchange data over the network.
%% This could be over only the local network or extended to the wider Internet.
%% As discussed previously, the communication between IoT services and the network may be in the form of a web service endpoint, an API, or in more pathological forms, such as files and HTTP queries.
%% Therefore, while we consider attackers who can only reach the target IoT device over the network, unlike \karonte or \satc, we do not limit \emph{how} attackers reach the target device.
%% We focus on any binary that could potentially process external data in a harmful manner, be it through sockets, arguments passed through command lines, or environment and NVRAM variables.


\smallskip
\noindent
\textbf{Bug reports and triggerable vulnerabilities.}
Because \mango runs on every binary in a given firmware, it identifies legitimate bugs in the current firmware regardless of whether the bugs are triggerable.
For example, \mango can identify a vulnerable function that is unused or in an executable on the file system that an attacker cannot currently execute.
While finding these bugs is essential, in this paper, we give a higher priority to finding \emph{triggerable} vulnerabilities~\footnote{When we conduct responsible disclosure for vulnerabilities to vendors, they would only accept vulnerability reports that come with Proof-of-Concept exploits (PoCs) that trigger the vulnerability and proactively reject vulnerability reports without PoCs.}.
We coin these triggerable vulnerabilities \emph{True PoCable vulnerabilities}, or \emph{\trupoc}s in short.
We implement a filtering system (described in Section \ref{sec:alert-ranking}) to prioritize bug reports that are more likely to be reachable by an attacker, thereby separating triggerable vulnerabilities from ordinary bug reports.
%% We argue that (1) determining the reachability of a vulnerability in real-world binaries is an undecidable problem in general: It is possible that we simply could not find the input to trigger the vulnerability, but a more skillful attacker might, and (2) modern embedded devices frequently receive updates throughout their lifetime, which means a presently unreachable vulnerability may become exploitable in the future.

